\cleardoublepage

\section{Introduction}

\subsection{Background}

Since mid-2024, a special phenomenon has been observed again and again in the LLM and Linguistics community. Multiple frontier LLMs including GPT-4, Claude 3.5, and Llama 3 systematically answer "which is bigger, 9.11 or 9.9" a seemingly very absurd wrong answer: 9.11 is bigger\cite{singh2024, eval16x2025}. This error is very common. It is not an occasional phenomenon of some special model under some special setting. It can be stably reproduced with different temperature parameters, different prompt templates, and even different languages. On examples where the numerical gap is smaller (such as 8.7 vs 8.12, 3.11 vs 3.9), it shows predictable transferability. What is even more thought-provoking is that this error is not unfamiliar. As early as twenty years ago, Desmet et al.\cite{desmet2010} had recorded almost exactly the same phenomenon in upper primary school and junior high school students. A considerable proportion of children would insist that 0.25 is bigger than 0.8, because 25 is bigger than 8. There is a certain similarity between the composition of LLMs and the human brain: they are both made up of neurons. However, their actual working mechanisms have very big differences. LLMs are artificially designed neural networks made up of the transformer architecture, and are trained on large corpora to make them gain rich natural language and world knowledge. The human brain, on the other hand, gradually acquires language and knowledge through restructured feedback with the real world. The core question this study aims to answer is: is this coincidence a surface similarity, or is it a mechanistic isomorphism? We hope that by studying and comparing the differences between humans and LLMs on number bias, we can produce valuable conclusions for the acquisition mechanism of human language and the development of LLMs.

\subsection{Natural Number Bias in Humans}

In math cognition research, the number errors made by the children above are grouped under the topic of natural number bias (NNB; also called whole number bias, WNB)\cite{alibali2015, alonsodaz2018}. NNB describes a class of systematic failures: when learners face rational numbers (decimals, fractions, percentages), they unknowingly carry over rules that only apply to natural numbers into the new number domain. A few clear examples are: discreteness, the successor relation, and the more digits means bigger error. To explain this problem, several theories have been proposed over the years to model it. Roughly, they can be grouped from phenomenon to mechanism, from description to explanation, into five main paths.

The first path is the early rule induction model. The earliest work that systematically recorded decimal comparison errors goes back to Brueckner\cite{brueckner1928}, who studied difficulties primary school students had with decimals. But the work that actually shaped the research framework that followed was Sackur-Grisvard and Leonard\cite{sackurgrisvard1985}, who studied how French children in grades 4 to 7 ordered decimals. Through qualitative analysis, the two researchers grouped the children's errors into three hidden rules. Rule 1 treats the decimal part as a whole number and says the decimal with the smaller reading is larger (for example, thinking 12.4 < 12.17 because 4 < 17). Rule 2 says the decimal with more digits in the decimal part is smaller (for example, judging 12.94 to be less than 12.7, which reflects an over-generalized idea that the hundredths place is smaller than the tenths place). Rule 3 treats a leading zero in the decimal part as a signal, and gives the right answer when this rule applies, but falls back to Rule 1 when it does not. Nesher and Peled\cite{nesher1986} repeated the test with Israeli students in grades 7 to 9, and for the first time allowed "equal" as an answer choice. They found that the use of Rule 1 and Rule 3 went down with grade level, while the use of Rule 2 went up in grade 8 and then went down again. This was the first empirical evidence that NNB is structured, can be classified, and changes with development. Resnick et al.\cite{resnick1989} further made the conceptual sources of these rules clearer: Rule 1 comes from negative transfer of natural number knowledge, and Rule 2 comes from over-generalizing the fraction rule that a bigger denominator means a smaller value. Work in this period was mostly descriptive and had not yet formed a unified explanation of the underlying mechanism. But it set up a research method that is still used today, which is to design "congruent vs. incongruent" conditions to separate surface features from real numerical value.

The second path is conceptual change theory. This framework comes from Kuhn\cite{kuhn1962} on scientific revolutions, was brought into science education by Posner, Strike, Hewson and Gertzog\cite{posner1982}, then developed by Vosniadou\cite{vosniadou1994} into a more general theory of cognitive development, and finally applied to math learning by Vamvakoussi and Vosniadou\cite{vamvakoussi2004} and Stafylidou and Vosniadou\cite{stafylidou2004}. Its core claims are three. First, learners come in with prior concepts. When new knowledge is introduced, these concepts are not simply overwritten, but compete with the new ones over a long time. Second, conceptual change is slow and gradual, and cannot be done in one teaching session. Third, before the change is finished, learners go through synthetic models, which are middle states where they force the old and new rules together. Desmet, Grégoire and Mussolin\cite{desmet2010} gave key support for this theory in a study of 561 Belgian children in grades 3 to 6. Using a fine-grained design with 80 decimal pairs across 8 condition types, they were the first to statistically separate the effect of length from the effect of value. They found that children in grades 3 and 4 got only 4.1\% correct on items where length and value were both incongruent (such as 0.1 vs 0.02), and most errors (95.3\%) were picking the longer number. In grade 5 a reverse bias appeared, where children thought the shorter decimal was larger. This is the statistical version of what Sackur-Grisvard called Rule 2. Only in grade 6 did accuracy reach a stable level of 97\% or more. A cluster analysis further found 5 types of learners, from Cluster 1 who only relied on numerical reading, to Cluster 4 who were close to expert level but still got stuck on trailing zeros, to Cluster 5 who had fully mastered the topic. This matched closely with the middle synthetic models that Vosniadou predicted. The contribution of this path is that it puts NNB as a long-term competition between two systems of representation (natural numbers vs. rational numbers), rather than just missing knowledge or carelessness.

The third path is misconception taxonomy, which can be seen as the engineering extension of conceptual change theory. Isotani, McLaren and Altman\cite{isotani2010} put together more than 80 years of work, about 40 papers (Brueckner 1928, Resnick 1989, Stacey et al. 2001, etc.), and grouped decimal errors into six classes named with mnemonics. Megz (Longer-is-Larger, like 0.25 > 0.7, comes from carrying over whole number rules). Segz (Shorter-is-Larger, like 0.2 > 0.25, comes from over-generalizing fraction knowledge). Negz (treating decimals less than 1 as if they were negative). Regz (treating decimals as whole numbers entirely). Ovegz (a basic misunderstanding of place value, like treating 0.67 as 67 tenths). Zegz (wrong rules about the position of zero). This taxonomy not only sorted out findings that had been spread across many papers, but was also built into the smart tutoring system MathTutor as a diagnostic dimension of the student user model. It supported the teaching method of erroneous examples (ErrEx), based on Ohlsson's\cite{ohlsson1996} theory of learning from execution errors, by having students compare correct and wrong solutions and explain the errors themselves to fix their concepts. The methodological value of this path is that it turns the descriptive concepts from developmental psychology into objects that can be computed, diagnosed, and acted on. This is exactly the most important infrastructure needed when the NNB framework is later applied to LLMs.

The fourth path is dual-processing theory. It tries to provide a real-time mechanism, at the level of processing, for the long-term coexistence of two representations described by conceptual change theory. The theory divides cognition into the intuitive system S1 (fast, automatic, parallel) and the analytic system S2 (slow, effortful, serial). NNB is explained as: S1 automatically activates a magnitude representation based on the size of natural numbers, and S2 fails to step in or correct it in time\cite{vamvakoussi2012}. This framework predicts reaction time patterns in magnitude comparison tasks fairly well: when the magnitude distance between two numbers is large, S1 directly gives the right answer; when the distance is small, S2 has to step in and use an effortful componential strategy, and reaction time goes up. But Alibali and Sidney\cite{alibali2015} point out in their review that this framework has limits. In a task where students have to judge algebraic expressions (for example, whether x < x+2 is true for all x), Van Hoof et al.\cite{van2015} found that students almost always plug in whole number values to test, and this process is clearly an effortful, explicit strategy that cannot be put under S1's automatic activation. In other words, NNB can also happen during analytic reasoning, and a pure two-system contrast cannot cover all the phenomena.

The fifth path is the dynamic strategy choice framework proposed by Alibali and Sidney\cite{alibali2015}, which is an integrative extension of the explanations above. This framework borrows from Siegler and Shipley's\cite{siegler1995} strategy choice model for arithmetic learning. It claims that when solving any rational number problem, a person has a strategy repertoire that includes both intuition-based magnitude retrieval strategies and various effortful componential strategies. The actual choice is decided by the current task affordances, the make-up of the person's strategy repertoire, and the recent history of solving similar problems. The key contribution of this framework is that it pushes the source of NNB further down to the representation level: certain strategy choices lead to NNB because the learner's mental representations of natural number magnitude and operational patterns are formed earlier, activated more often, with greater strength and more accuracy than those for rational numbers\cite{dehaene2009, siegler2011}. When a task pushes the learner to access a representation, but the rational number representation is not strong enough to take over the decision, the natural number representation steps in, and NNB shows up. The advantage of this framework is that it explains variability in two directions at once: why the same person gets different results across trials (task affordances change), and why people of different ages or experience show systematic differences in groups (strategy repertoire and representation strength develop over time).

The most recent path is the intrinsic whole number bias hypothesis proposed and tested through Bayesian modeling by Alonso-Díaz, Piantadosi, Hayden and Cantlon\cite{alonsodaz2018}. It stands against the strategic bias hypothesis, supported by Schneider \& Siegler\cite{schneider2010} and Fazio, DeWolf, \& Siegler\cite{fazio2016}. The strategic hypothesis says that when comparing fractions, people use a holistic ratio strategy by default, and only fall back on whole number comparison when the denominators are too close and the holistic ratio is hard to extract. In other words, NNB is a fall-back when this strategy fails. The intrinsic hypothesis instead says that in every ratio comparison, people compute both the holistic ratio and the component whole numbers at the same time, and use a weighted combination of the two to decide. So whole number interference is always there, and only its weight changes with the task. Alonso-Díaz et al.\cite{alonsodaz2018} ran a fast 2AFC probability comparison task with N=21 (816 trials), and put 5 competing explanations (denominator neglect, whole number weighting, holistic ratio, strategic, and intrinsic) into one Bayesian framework with different parameter constraints. Through MCMC sampling (Stan, 8 chains × 10,000 iters) and WAIC model comparison, they found that the intrinsic model was best at both the group and individual level. The 95\% credible intervals of decision weights showed that both the numerator weight (winners, [0.045, 0.099]) and the holistic ratio weight (ratios, [3.040, 4.263]) were clearly different from zero. This result was strengthened by an earlier finding (Exp 1, N=284): even when subjects correctly reported in an explicit question that the two urns had equal probability (error < 5\%), 66\% still preferred the urn with more balls (p = .001). In other words, WNB is not because people don't know. It's that even when they know, they still cannot escape it. The theoretical meaning of this path is that it pulls NNB from "not enough education or missing concepts" further down to the level of the inner structure of the core cognitive system. Together with non-verbal ratio cognition evidence from infants and primates\cite{denison2010, mccrink2007, rakoczy2014}, it suggests that whole number interference may come from an early-developing, cross-species non-verbal number sense system, rather than being just a product of culture or teaching.

\subsection{Natural Number Bias in LLMs and Interpretability}

On the LLM side, looking at the overall benchmark performance, frontier models have long stopped struggling with traditional math problems. On datasets like GSM8K (elementary school word problems) and MATH (competition-level problems), mainstream models' accuracy has generally passed 85\%\cite{cobbe2021, wallace2019}, and some models even get close to human expert level. But the truly strange phenomenon shows up on the most basic numerical comparison tasks. On questions like "which is bigger, 9.11 or 9.9" or "is 3.8 greater than 3.11", which need almost no reasoning chain and only involve reading decimal places and comparing sizes, the models often make mistakes. Several studies have already noticed this: Xie\cite{xie2602} and 16x Eval\cite{eval16x2025}, among others, have documented from different angles the error patterns of models on simple number bias problems, including confusion over decimal places. Since this phenomenon is very similar to the NNB errors made by humans, we should analyze them together.

However, since an LLM is a neural network with a huge number of parameters, its internal state is very hard to track and explain. But recent progress in the mechanistic interpretability field gives us tools to look at this problem again that we did not have before. Linear probing trains a simple linear classifier on the model's hidden states, and can tell whether a certain concept is encoded in a linearly separable way in the activations of some layer. Activation patching replaces the activation values at specific positions between two forward passes, locating which components and which positions are causally responsible for a certain judgment. With these methods, researchers have already gotten relatively clear internal correspondences for human-understandable concepts like numerical size, spatial direction, and true/false judgment\cite{meng2022, huang2025}. This means that, for the question of "what the model actually read, what it compared, and at which step it went off track" at the moment it judges whether 9.11 or 9.9 is bigger, we can directly open up the model and observe the real trajectory of its internal state when it makes the mistake.
